\documentclass{article}
% ---------------------------------------------------------------
% Póster de revisión de literatura (A0 vertical)
% Finanzas I (055D), UNCP, 2026-II. Compilar con pdfLaTeX.
% ---------------------------------------------------------------
\usepackage[paperwidth=841mm,paperheight=1189mm,margin=0mm]{geometry}
\usepackage[utf8]{inputenc}
\usepackage[T1]{fontenc}
\IfFileExists{spanish.ldf}{\usepackage[spanish,es-noshorthands,es-nodecimaldot]{babel}}{}
\usepackage{helvet}
\renewcommand{\familydefault}{\sfdefault}
\usepackage[letterspace=120]{microtype}
\usepackage{anyfontsize}
\usepackage{graphicx}
\usepackage{xcolor}
\usepackage[most]{tcolorbox}
\usepackage{enumitem}
\usepackage{multicol}
\usepackage{xurl}
\usepackage[hidelinks]{hyperref}
\usepackage{qrcode}
\urlstyle{same}
\pagestyle{empty}
\setlength{\parindent}{0pt}
\hyphenpenalty=10000 \exhyphenpenalty=10000

% ---------------- Paleta ----------------
\definecolor{navy}{HTML}{8E0C11}
\definecolor{navyLight}{HTML}{B8121A}
\definecolor{gold}{HTML}{00A3E0}
\definecolor{acento}{HTML}{A6E1F7}
\definecolor{goldSoft}{HTML}{E8F6FC}
\definecolor{teal}{HTML}{353435}
\definecolor{tealSoft}{HTML}{F1EEEE}
\definecolor{ink}{HTML}{231F20}
\definecolor{muted}{HTML}{5E5757}
\definecolor{line}{HTML}{E5D7D7}
\definecolor{page}{HTML}{F6F1F1}
\definecolor{mist}{HTML}{F7D3D5}
\pagecolor{page}

% ---------------- Tamaños (ajustados por póster) ----------------
\newcommand{\bodysize}{\fontsize{27}{35.6}\selectfont}
\newcommand{\cardtitlesize}{\fontsize{31.9}{38.3}\selectfont}
\newcommand{\smallsize}{\fontsize{22.1}{28.7}\selectfont}

% ---------------- Estilos de tarjeta ----------------
\tcbset{
  card/.style={enhanced, sharp corners, colback=white, colframe=line, boxrule=0.9mm,
    left=14mm, right=14mm, top=12mm, bottom=12mm, fontupper=\bodysize\color{ink}\raggedright,
    overlay={\fill[#1] (frame.north west) rectangle ([yshift=-5mm]frame.north east);}},
  card/.default=navy,
}
\newcommand{\cardtitle}[2][navy]{{\color{#1}\cardtitlesize\bfseries\textls[110]{\MakeUppercase{#2}}}\par\vspace{7mm}}
\newcommand{\bul}{\raisebox{0.18em}{\color{gold}\rule{0.5em}{0.5em}}}
\setlist[itemize]{label=\bul, leftmargin=1.25em, labelsep=0.6em, itemsep=0.55em, topsep=0pt, parsep=0pt}
\newcommand{\pill}[3]{\tcbox[on line, enhanced, arc=7mm, boxrule=0.8mm, colback=#1, colframe=#2,
  left=7mm, right=7mm, top=3mm, bottom=3mm, fontupper=\fontsize{26}{30}\selectfont\bfseries]{#3}}

\begin{document}

% ======================= ENCABEZADO =======================
\begin{tcolorbox}[enhanced, sharp corners, colback=navy, boxrule=0pt, width=\paperwidth,
  left=34mm, right=34mm, top=20mm, bottom=20mm,
  overlay={\fill[gold] (frame.south west) rectangle ([yshift=5mm]frame.south east);
           \fill[navyLight] ([xshift=-160mm]frame.north east) -- (frame.north east) -- ([yshift=-160mm]frame.north east) -- cycle;}]
\begin{minipage}[c]{92mm}\centering\includegraphics[width=\linewidth]{Escudo_UNCP.png}\end{minipage}\hfill
\begin{minipage}[c]{555mm}\centering
{\color{acento}\fontsize{25}{30}\selectfont\bfseries\textls[160]{REVISIÓN DE LITERATURA \quad|\quad FINANZAS I \quad|\quad TEMA 61}\par}
\vspace{10mm}
{\color{white}\fontsize{58}{65.5}\selectfont\bfseries Beta estimation in the European network regulation context: what matters, what doesn't, and what is indispensable\par}
\vspace{8mm}
{\color{mist}\fontsize{34}{41.5}\selectfont\itshape Estimación de la beta en el contexto de la regulación europea de redes: qué importa, qué no y qué es indispensable\par}
\vspace{10mm}
{\color{white}\fontsize{29}{35}\selectfont\bfseries Dmitry Bazhutov, André Betzer y Richard Stehle\par}
\vspace{9mm}
\pill{white}{white}{\color{navy}Q2}\hspace{6mm}\pill{navy}{white}{\color{white}\mdseries \textit{Financial Markets and Portfolio Management}, 37(3), 239-275 (2023)}
\end{minipage}\hfill
\begin{minipage}[c]{92mm}\centering\includegraphics[width=\linewidth]{Logo_FE.png}\end{minipage}
\end{tcolorbox}
\vspace{-1mm}
\begin{tcolorbox}[enhanced, sharp corners, colback=goldSoft, boxrule=0pt, width=\paperwidth,
  left=34mm, right=34mm, top=8mm, bottom=8mm, halign=center]
{\color{navy}\fontsize{27}{36}\selectfont \textbf{Póster elaborado por:} Meglinho Sneyder Rojas Poma \enspace\textbar\enspace Código 2024200522B\par \fontsize{24}{32}\selectfont Escuela Profesional de Economía, Universidad Nacional del Centro del Perú \enspace\textbar\enspace Finanzas I (055D), 2026-II\par}
\end{tcolorbox}

\vfill
% ======================= CUERPO =======================
\hspace*{34mm}\begin{minipage}{773mm}
\begin{tcolorbox}[enhanced, sharp corners, colback=navy, boxrule=0pt, left=16mm, right=16mm, top=12mm, bottom=11mm,
  fontupper=\bodysize\raggedright, overlay={\fill[gold] (frame.north west) rectangle ([yshift=-5mm]frame.north east);}]
\begin{minipage}[t]{174.8mm}\raggedright{\color{acento}\smallsize\bfseries\textls[110]{\MakeUppercase{Mercado}}}\par\vspace{2mm}{\color{white}Operadores de redes regulados en Europa}\end{minipage}\hfill\begin{minipage}[t]{174.8mm}\raggedright{\color{acento}\smallsize\bfseries\textls[110]{\MakeUppercase{Muestra}}}\par\vspace{2mm}{\color{white}Acciones individuales y grupos pares}\end{minipage}\hfill\begin{minipage}[t]{174.8mm}\raggedright{\color{acento}\smallsize\bfseries\textls[110]{\MakeUppercase{Periodo}}}\par\vspace{2mm}{\color{white}Hasta 2020, con el inicio del COVID-19}\end{minipage}\hfill\begin{minipage}[t]{174.8mm}\raggedright{\color{acento}\smallsize\bfseries\textls[110]{\MakeUppercase{Método}}}\par\vspace{2mm}{\color{white}Betas por MCO con distintas ventanas}\end{minipage}
\end{tcolorbox}
\end{minipage}

\vfill
\hspace*{34mm}\begin{minipage}{773mm}
\begin{tcbraster}[raster columns=3, raster equal height, raster column skip=16mm]
\begin{tcolorbox}[card]\cardtitle{Contexto}En Europa, las empresas de redes de gas, electricidad y telecomunicaciones están reguladas, y el regulador fija sus tarifas a partir del costo de capital. Ese costo se calcula con el CAPM, por lo que la beta estimada tiene consecuencias económicas directas.\end{tcolorbox}
\begin{tcolorbox}[card]\cardtitle{Problema}La beta no se observa: se estima, y el resultado depende de decisiones como la ventana de tiempo, la frecuencia de los retornos y el índice de mercado elegido. Pocas investigaciones revisan esas decisiones en detalle para un grupo pequeño de empresas.\end{tcolorbox}
\begin{tcolorbox}[card]\cardtitle{Objetivo}Evaluar paso a paso cómo se estima la beta, desde la selección y limpieza de datos hasta los ajustes finales, usando acciones individuales y grupos pares de operadores de redes europeos.\end{tcolorbox}
\end{tcbraster}
\end{minipage}

\vfill
\hspace*{34mm}\begin{minipage}{773mm}
\begin{tcolorbox}[card]\cardtitle{Metodología}
\begin{minipage}[t]{0.48\linewidth}
\begin{itemize}
\item Ventanas de uno, tres y cinco años.
\item Retornos diarios, semanales y mensuales.
\item Índices de mercado locales, europeos y globales.
\end{itemize}
\end{minipage}\hfill
\begin{minipage}[t]{0.48\linewidth}
\begin{itemize}
\item Comparación con los ajustes de Vasicek y Blume.
\item Se analizan los efectos de la crisis financiera de 2008 y del inicio del COVID-19 en 2020 sobre las series de betas.
\item Se incluyen pares de electricidad, gas, telecomunicaciones y agua de Europa, y ferrocarriles de EE. UU. como comparación.
\end{itemize}
\end{minipage}
\end{tcolorbox}
\end{minipage}

\vfill
\hspace*{34mm}\begin{minipage}{773mm}
\begin{tcolorbox}[card]\cardtitle{Figura clave del estudio}
\begin{minipage}[c]{375.5mm}\centering
\fcolorbox{line}{white}{\includegraphics[width=\dimexpr375.5mm-2\fboxsep-2\fboxrule\relax]{figura.png}}
\end{minipage}\hfill
\begin{minipage}[c]{\dimexpr\linewidth-375.5mm-18mm\relax}\raggedright
{\color{navy}\bfseries Figura 2.}\enspace{\itshape Betas semanales a cinco años de Enagás, S.A. con el índice nacional (IBEX 35), europeo (Stoxx Europe 600) y global (MSCI World), de diciembre de 2010 a junio de 2020.}\par\vspace{8mm}
{\color{teal}\smallsize\bfseries\textls[110]{CÓMO LEERLA}}\par\vspace{4mm}
\begin{itemize}
\item La beta media es 0.60 con el IBEX 35, 0.72 con el Stoxx Europe 600 y 0.68 con el MSCI World.
\item Con el índice europeo la beta llegó a 0.91, mientras que con el nacional no pasó de 0.67.
\item La misma empresa, en el mismo periodo, tiene betas distintas según el índice de mercado elegido.
\end{itemize}
\vspace{8mm}
{\color{muted}\smallsize Figura tomada de Bazhutov et al. (2023).}
\end{minipage}
\end{tcolorbox}
\end{minipage}

\vfill
\hspace*{34mm}\begin{minipage}{773mm}
\begin{tcolorbox}[card=teal]\cardtitle[teal]{Resultados}
\begin{tcolorbox}[enhanced, sharp corners, colback=tealSoft, boxrule=0pt, left=14mm, right=14mm, top=10mm, bottom=11mm,
  borderline west={4mm}{0mm}{teal}, fontupper=\bodysize\color{ink}\raggedright]
{\color{teal}\smallsize\bfseries\textls[110]{HALLAZGO CENTRAL}}\par\vspace{4mm}
{\bfseries Las betas a tres y, sobre todo, a cinco años son mucho más estables que las de un año, y las betas semanales o diarias funcionan mejor que las mensuales.}
\end{tcolorbox}
\vspace{11mm}
\begin{minipage}[t]{0.48\linewidth}
\begin{itemize}
\item Aparecen subidas y bajadas bruscas de la beta que duran poco y pueden distorsionar la predicción.
\item Elegir un índice local, europeo o global puede cambiar bastante la beta de algunas empresas.
\item Los ajustes de Vasicek y Blume hacia 1 producen betas futuras demasiado altas.
\end{itemize}
\end{minipage}\hfill
\begin{minipage}[t]{0.48\linewidth}
\begin{itemize}
\item Da igual usar retornos simples o logarítmicos, o retornos totales o en exceso.
\item Los retornos semanales o diarios parecen mejores que los mensuales.
\item Cuando los datos de una crisis salen de la ventana de estimación, aparecen caídas bruscas de la beta.
\end{itemize}
\end{minipage}
\end{tcolorbox}
\end{minipage}

\vfill
\hspace*{34mm}\begin{minipage}{773mm}
\begin{tcbraster}[raster columns=3, raster equal height, raster column skip=16mm]
\begin{tcolorbox}[card, raster multicolumn=2]\cardtitle{Conclusiones}Se recomienda usar retornos totales, expresar todo en la misma moneda, tratar los días sin negociación y evitar las betas de un año en la regulación, por su fuerte variación.\end{tcolorbox}
\begin{tcolorbox}[card=gold, colback=goldSoft]\cardtitle{Aporte al tema 61}Funciona como guía práctica para el tema 61: justifica convertir todos los precios a una sola moneda, filtrar los días sin negociación y comprobar la beta con otro índice de mercado.\end{tcolorbox}
\end{tcbraster}
\end{minipage}

\vfill
% ======================= PIE =======================
\begin{tcolorbox}[enhanced, sharp corners, colback=navy, boxrule=0pt, width=\paperwidth,
  left=34mm, right=34mm, top=14mm, bottom=11mm,
  overlay={\fill[gold] (frame.north west) rectangle ([yshift=-5mm]frame.north east);}]
\begin{minipage}[c]{245mm}\raggedright
{\color{acento}\smallsize\bfseries\textls[110]{FUENTE DE DATOS DEL ESTUDIO}}\par\vspace{5mm}
{\color{white}\smallsize Cotizaciones de operadores de redes europeos e índices Stoxx Europe 600 y Stoxx Europe TMI.}
\end{minipage}\hfill
\begin{minipage}[c]{385mm}\raggedright
{\color{acento}\smallsize\bfseries\textls[110]{REFERENCIA}}\par\vspace{5mm}
{\color{white}\smallsize\hangindent=1.2em\hangafter=1 Bazhutov, D., Betzer, A., y Stehle, R. (2023). Beta estimation in the European network regulation context: What matters, what doesn't, and what is indispensable. \textit{Financial Markets and Portfolio Management}, \textit{37}(3), 239-275. \url{https://doi.org/10.1007/s11408-023-00428-z}\par}
\end{minipage}\hfill
\begin{minipage}[c]{100mm}\centering
\tcbox[enhanced, sharp corners, colback=white, boxrule=0pt, left=3mm, right=3mm, top=3mm, bottom=3mm]{\qrcode[hyperlink, height=46mm]{https://doi.org/10.1007/s11408-023-00428-z}}\par\vspace{4mm}
{\color{acento}\fontsize{17}{21}\selectfont\bfseries\textls[100]{ESCANEA Y LEE EL ARTÍCULO}}
\end{minipage}
\end{tcolorbox}
\end{document}
